\par\medskip
\noindent\textbf{4.14.}
Cerramos esta sección construyendo explícitamente una deformación
formal versal en un caso sencillo. Cuando $X_0$ es afín y una
intersección completa, tal deformación se escribe de inmediato. Sea
\[
A_0=k[x_1,\ldots,x_n]/(f_1,\ldots,f_r)=S/J,
\]
donde $S=k[x_1,\ldots,x_n]$, $J=(f_1,\ldots,f_r)$ y
$f_1,\ldots,f_r$ es una sucesión $S$-regular. Tenemos una sucesión
exacta
\[
\cdots\longrightarrow\operatorname{Der}_k(S,A_0)
\longrightarrow\operatorname{Hom}_{A_0}(J/J^2,A_0)
\longrightarrow D^1(A_0|k,A_0)\longrightarrow0.
\]
El $A_0$-módulo $J/J^2$ es libre de rango $r$, generado por las
clases de $f_1,\ldots,f_r$, porque estos elementos forman una
sucesión $S$-regular.

Supongamos ahora que $D^1(A_0|k,A_0)$ es de dimensión finita sobre
$k$; por ejemplo, supongamos que $A_0$ solo tiene puntos no lisos
aislados. Pongamos
\[
d=\dim_kD^1(A_0|k,A_0),
\]
y elijamos aplicaciones
\[
p_i:J/J^2\longrightarrow A_0\qquad(1\leq i\leq d)
\]
cuyas clases formen una $k$-base de $D^1(A_0|k,A_0)$. Para
$1\leq i\leq d$ y $1\leq j\leq r$, elijamos
$P_{ij}\in\Lambda[x_1,\ldots,x_n]$ tal que
\[
P_{ij}\equiv p_i(f_j)\pmod{\mathfrak m_\Lambda+J},
\]
donde $\mathfrak m_\Lambda$ es el ideal maximal de $\Lambda$. Pongamos
\[
R=\Lambda[[t_1,\ldots,t_d]],
\qquad
B=R[x_1,\ldots,x_n]/(G_1,\ldots,G_r),
\]
donde
\[
G_j=F_j-\sum_{i=1}^d t_iP_{ij},
\]
y $F_j\in\Lambda[x_1,\ldots,x_n]$ se obtiene de $f_j$ levantando sus
coeficientes de $k$ a $\Lambda$.

El álgebra $R[x_1,\ldots,x_n]$ es plana sobre $R$, y las imágenes de
$G_1,\ldots,G_r$ módulo $\mathfrak m_R$ forman la sucesión regular
$f_1,\ldots,f_r$ en
$k\otimes_RR[x_1,\ldots,x_n]$. Por tanto, $B$ es $R$-plana y
$k\otimes_RB=A_0$. Por construcción, la aplicación
\[
\operatorname{Hom}_{\Lambda\text{-alg}}(R,k[\varepsilon])
\longrightarrow D^1(A_0|k,A_0)
=[\underline M_{\operatorname{Spec}(A_0)}(k[\varepsilon])]
\]
es biyectiva. Así, si

%%%%%%%%%%  scan idx 122  |  folio 111  %%%%%%%%%%

\[
\operatorname{Spf}(B)
=\varinjlim_\nu
\operatorname{Spec}\bigl((R/\mathfrak m_R^{\nu+1})\otimes_RB\bigr),
\]
entonces $(R,\operatorname{Spf}(B))$ es un objeto de
$\widehat{\underline M}_{\operatorname{Spec}(A_0)}$. El funtor
\[
\underline h_R\longrightarrow\underline M_{\operatorname{Spec}(A_0)}
\]
que induce es liso por 1.6, puesto que $R$ es un anillo de series
formales sobre $\Lambda$. En otras palabras,
$(R,\operatorname{Spf}(B))$ es un objeto formal versal de
$\underline M_{\operatorname{Spec}(A_0)}$.

Por 4.10(b), si ponemos en cambio
\[
B=R[[x_1,\ldots,x_n]]/(G_1,\ldots,G_r),
\]
entonces $(R,\operatorname{Spf}(B))$ es una deformación formal versal
de $\operatorname{Spec}(A_0)$ para
\[
A_0=k[[x_1,\ldots,x_n]]/(f_1,\ldots,f_r).
\]

\medskip
\noindent\textbf{Observación 4.15.}
Sea $X_0$ un esquema de tipo finito sobre $k$ que solo tiene puntos
no lisos aislados, y sea $(R,\mathfrak X)$ una deformación formal
versal de $X_0$. La construcción precedente muestra que, cuando
$X_0$ es afín y localmente una intersección completa, $\mathfrak X$
es la formalización de un $R$-esquema de tipo finito. El mismo hecho
se extiende fácilmente al caso no afín siempre que
$H^2(X_0,\underline D^0_{X_0})=0$.
